\begin{apartats}

\apartat[2,5]{1,5}
Determina els valors de $a$ i $b$ perquè la funció següent sigui contínua a tots els
punts de $\mathbb{R}$.
\[
f(x)=\begin{cases}
  x^2-a & \si{x<-1},\\[3pt]
  bx+2  & \si{-1\le x\le 2},\\[3pt]
  ax+b  & \si{x>2}.
\end{cases}
\]

\begin{solucio}
Cada branca és contínua al seu tros, de manera que només cal imposar la continuïtat
als enganxaments.\\
En $x=-1$: $\lim_{x\to-1^-}f(x)=1-a$ i $f(-1)=-b+2$, d'on $1-a=-b+2$, és a dir $b=a+1$.\\
En $x=2$: $f(2)=2b+2$ i $\lim_{x\to2^+}f(x)=2a+b$, d'on $2b+2=2a+b$, és a dir $b=2a-2$.\\
Igualant: $a+1=2a-2\Rightarrow \boxed{a=3}$ i $\boxed{b=4}$.
\end{solucio}

\begin{nomesllarg}
\apartat{1}
Troba el valor de $k$ perquè la funció següent sigui contínua a tot $\mathbb{R}$.
\[
g(x)=\begin{cases}
  kx+1  & \si{x\le 2},\\[3pt]
  x^2-k & \si{x>2}.
\end{cases}
\]

\begin{solucio}
Les dues branques són polinòmiques: només cal estudiar $x=2$.\\
$g(2)=2k+1$ i $\lim_{x\to2^+}g(x)=4-k$. Cal $2k+1=4-k$, és a dir $3k=3$: $\boxed{k=1}$.
(Comprovació: $g(2)=3$ i $4-1=3$.)
\end{solucio}
\end{nomesllarg}

\end{apartats}
